\documentclass[11pt]{article}
\usepackage{amsmath}
\usepackage{amssymb}
\usepackage[margin=1in]{geometry}
\usepackage{hyperref}

\title{Methodology for Estimating Global Distributed Desalination Costs}
\author{Paraphrased from Kocher \& Menon, \emph{Energy Environ. Sci.}, 2023, 16, 4983--4993}
\date{}

\begin{document}
\maketitle

\section{Levelized Cost of Water for Coastal Desalination}

The baseline cost of producing freshwater via seawater reverse osmosis (RO) at
the coast is taken as approximately \$1 per m$^3$. This figure is derived from
data on operating RO plants and includes both capital expenditure (CAPEX) and
operating expenditure (OPEX), where OPEX incorporates the levelized cost of
electricity (LCOE) needed to drive the RO process, as well as the cost of
disposing of the resulting brine (typically via ocean discharge).

\section{Extending the Cost to Inland (Distributed) Locations}

Because desalination facilities are sited at the coast but a majority of the
global population lives inland, the analysis adds the cost of conveying the
treated water from the nearest coastline to each inland location. The nearest
coastal point for any given location is identified using a high-resolution
global coastline dataset (the Global Self-Consistent Hierarchical
High-Resolution Geography, or GSHHG, as implemented in MATLAB's mapping
toolbox).

The transport cost itself is levelized over the operating lifetime of the
conveyance infrastructure (a mix of pipes, canals, and tunnels), combining the
relevant capital and operating expenditures, so that it can be expressed on
the same \$ per m$^3$ basis as the desalination cost and added directly to it.

\subsection{Vertical and Horizontal Transport Costs}

Transport costs are broken into a vertical component (pumping water uphill in
elevation) and a horizontal component (moving water laterally overland),
based on a literature review of water conveyance economics. The baseline
(unit) costs adopted are:

\begin{align}
c_{\text{vert}} &\approx 5 \times 10^{-4}\ \$/\text{m}^3\text{ per m of vertical rise} \\
c_{\text{horiz}} &\approx 6 \times 10^{-4}\ \$/\text{m}^3\text{ per km of horizontal distance (canal basis)}
\end{align}

The horizontal figure above corresponds to open canals. Because real-world
conveyance often uses a mix of canals, tunnels (reported as roughly twice as
costly as canals), and pipes (reported as roughly 3.7$\times$ the cost of
canals), the study assumes an equal-weighted blend of the three
transport modes. This yields an effective horizontal transport cost of

\begin{equation}
c_{\text{horiz,eff}} \approx 13.58 \times 10^{-4}\ \$/\text{m}^3\text{ per km}.
\end{equation}

This blended assumption is motivated by real infrastructure examples (e.g.,
the Colorado River Aqueduct in the United States), which combine multiple
conveyance types over their length.

\subsection{Total Delivered Cost}

For a given inland location, the total levelized cost of delivered water is
therefore the sum of the coastal RO production cost and the levelized
transport cost associated with the vertical and horizontal distance to the
nearest coast:

\begin{equation}
\text{LCOW}_{\text{delivered}} = \text{LCOW}_{\text{RO}} \;+\; c_{\text{vert}} \cdot \Delta z \;+\; c_{\text{horiz,eff}} \cdot d
\end{equation}

where $\Delta z$ is the vertical elevation change (m) and $d$ is the
horizontal distance (km) from the nearest coastline to the location of
interest.

\subsection{Sensitivity to Transport Cost Uncertainty}

Because conveyance costs reported in the literature vary considerably, the
authors also test a high-cost scenario in which both the vertical and
horizontal unit transport costs are increased by an order of magnitude
(i.e., $5\times10^{-3}$ \$/m$^3$ per m vertical and
$13.58\times10^{-3}$ \$/m$^3$ per km horizontal). This scenario is used to
identify how much more expensive transport would need to be before
distributed atmospheric water harvesting becomes cost-competitive in
particular regions (e.g., remote arid locations such as the Sahara).

\section{Application at Global Scale}

To generate global cost maps, this delivered-cost calculation is repeated for
every cell of a gridded representation of the Earth's surface, using each
cell's distance and elevation relative to the nearest coast. The resulting
grid of delivered desalination costs is then compared, cell by cell, against
the levelized cost of atmospheric water harvesting technologies to determine
which approach is cheaper at each location, and the results are further
weighted by population and water-risk data to assess the fraction of the
water-stressed global population each approach could serve most
cost-effectively.

\end{document}
